\chapter{拓扑学}

拓扑学研究的是在连续变形下保持不变的空间性质。和度量空间相比，拓扑空间不一定预先给出距离函数；它只保留“哪些集合应当被看作开集”这一层结构。因此，拓扑学提供了一种比欧氏空间更抽象、也更灵活的语言，用来讨论连续性、极限、连通性、紧致性以及空间之间的等价关系。

从分析学的角度看，拓扑学可以理解为对“开集方法”的彻底抽象。只要知道开集是什么，就可以重新定义闭集、邻域、收敛、连续映射和紧致性。这样一来，许多原本依赖坐标或距离的定理，就能被放到更一般的空间中重新表述。

\section{从度量空间到拓扑空间}

\subsection{度量空间中的开集}

设 $(X,d)$ 是一个度量空间。对 $x\in X$ 和 $r>0$，开球定义为
\[
    B_r(x)=\{y\in X:d(x,y)<r\}.
\]
子集 $U\subseteq X$ 称为开集，如果对每个 $x\in U$，都存在 $r>0$，使得
\[
    B_r(x)\subseteq U.
\]
这个定义只关心点附近是否能放入一个足够小的开球。换句话说，开集表达的是一种“局部可活动空间”：在 $U$ 中的任一点附近，都有一小块区域仍然完全留在 $U$ 内。

\begin{example}
在通常度量下，区间 $(0,1)$ 是 $\mathbb{R}$ 中的开集；而 $[0,1]$ 不是开集，因为端点 $0,1$ 附近的任意小开区间都会伸出 $[0,1]$ 之外。
\end{example}

\begin{proposition}[度量空间中连续性的拓扑形式]
设 $(X,d_X)$ 和 $(Y,d_Y)$ 是度量空间。映射 $f:X\to Y$ 连续，当且仅当对 $Y$ 中任意开集 $V$，原像
\[
    f^{-1}(V)=\{x\in X:f(x)\in V\}
\]
都是 $X$ 中的开集。
\end{proposition}

这个命题是通往一般拓扑空间的关键。它表明，连续性可以完全用开集的原像来描述，而不必显式使用距离。因此，如果某个集合 $X$ 上已经指定了哪些子集是开集，那么就可以谈论连续性。

\section{拓扑空间}

\subsection{拓扑的定义}

\begin{definition}[拓扑]
设 $X$ 是一个集合。$X$ 上的一个\textbf{拓扑}是 $\mathcal{T}\subseteq\mathcal{P}(X)$，满足：
\begin{enumerate}
    \item $\varnothing\in\mathcal{T}$ 且 $X\in\mathcal{T}$；
    \item 任意多个 $\mathcal{T}$ 中集合的并仍在 $\mathcal{T}$ 中；
    \item 有限多个 $\mathcal{T}$ 中集合的交仍在 $\mathcal{T}$ 中。
\end{enumerate}
二元组 $(X,\mathcal{T})$ 称为\textbf{拓扑空间}。$\mathcal{T}$ 中的元素称为开集。
\end{definition}

这个定义把度量空间中开集的基本性质抽象出来。注意第二条允许任意并，而第三条只要求有限交。原因是任意多个开集的交不一定开，例如
\[
    \bigcap_{n=1}^{\infty}\left(-\frac1n,\frac1n\right)=\{0\}
\]
在 $\mathbb{R}$ 的通常拓扑中就不是开集。

\begin{example}
若 $(X,d)$ 是度量空间，把所有由 $d$ 定义的开集收集起来，就得到 $X$ 上的一个拓扑，称为由度量 $d$ 诱导的拓扑。
\end{example}

\subsection{基本例子}

\begin{example}[离散拓扑]
对任意集合 $X$，令
\[
    \mathcal{T}=\mathcal{P}(X).
\]
也就是说，$X$ 的每个子集都是开集。这个拓扑称为\textbf{离散拓扑}。在离散拓扑中，每个点都是孤立的。
\end{example}

\begin{example}[平凡拓扑]
令
\[
    \mathcal{T}=\{\varnothing,X\}.
\]
这个拓扑称为\textbf{平凡拓扑}或不可分拓扑。它只承认空集和全集是开集，因此区分点的能力极弱。
\end{example}

\begin{example}[$\mathbb{R}^n$ 上的通常拓扑]
$\mathbb{R}^n$ 上由欧氏距离
\[
    d(x,y)=\sqrt{\sum_{i=1}^n(x_i-y_i)^2}
\]
诱导出的拓扑称为通常拓扑。分析学中讨论极限、连续和微分时，默认使用的基本上都是这个拓扑。
\end{example}

\begin{example}[余有限拓扑]
设 $X$ 是无限集合。令 $\mathcal{T}$ 由 $\varnothing$ 以及所有补集有限的集合组成，即
\[
    \mathcal{T}=\{\varnothing\}\cup\{U\subseteq X:X\setminus U\text{ 是有限集}\}.
\]
这给出一个拓扑，称为\textbf{余有限拓扑}。它说明，拓扑空间可以远比度量空间怪异。
\end{example}

\begin{example}[特定点拓扑]
固定一点 $p\in X$。令开集为 $\varnothing$ 以及所有包含 $p$ 的子集。也就是说，
\[
    \mathcal{T}=\{\varnothing\}\cup\{U\subseteq X:p\in U\}.
\]
这是一个拓扑，称为特定点拓扑。
\end{example}

\subsection{拓扑的比较}

\begin{definition}[更细与更粗]
设 $\mathcal{T}_1$ 和 $\mathcal{T}_2$ 是同一集合 $X$ 上的两个拓扑。如果
\[
    \mathcal{T}_1\subseteq\mathcal{T}_2,
\]
则称 $\mathcal{T}_2$ 比 $\mathcal{T}_1$ \textbf{更细}，也称 $\mathcal{T}_1$ 比 $\mathcal{T}_2$ \textbf{更粗}。
\end{definition}

更细的拓扑有更多开集，因此可以表达更强的局部区分能力。离散拓扑是最细的拓扑，平凡拓扑是最粗的拓扑。

\begin{example}
在任意集合 $X$ 上，都有
\[
    \{\varnothing,X\}\subseteq\mathcal{T}\subseteq\mathcal{P}(X),
\]
其中 $\mathcal{T}$ 是 $X$ 上任意拓扑。
\end{example}

\begin{remark}
拓扑越细，映射从该空间出发时越容易连续；拓扑越粗，映射进入该空间时越容易连续。这是因为连续性要求开集原像为开集。
\end{remark}

\section{基与子基}

\subsection{拓扑的基}

在实际使用中，直接列出所有开集通常很困难。更常见的做法是给出一族较小的基本开集，再用它们生成全部开集。

\begin{definition}[基]
设 $X$ 是集合。子集族 $\mathcal{B}\subseteq\mathcal{P}(X)$ 称为 $X$ 上某个拓扑的\textbf{基}，如果：
\begin{enumerate}
    \item 对每个 $x\in X$，存在 $B\in\mathcal{B}$ 使得 $x\in B$；
    \item 若 $x\in B_1\cap B_2$，其中 $B_1,B_2\in\mathcal{B}$，则存在 $B_3\in\mathcal{B}$，使得
    \[
        x\in B_3\subseteq B_1\cap B_2.
    \]
\end{enumerate}
由 $\mathcal{B}$ 生成的拓扑定义为：$U\subseteq X$ 开，当且仅当对每个 $x\in U$，存在 $B\in\mathcal{B}$ 使得 $x\in B\subseteq U$。
\end{definition}

\begin{example}
$\mathbb{R}$ 上所有开区间 $(a,b)$ 构成通常拓扑的一组基。事实上，每个通常意义下的开集都可以写成若干开区间的并。
\end{example}

\begin{proposition}
由一组基 $\mathcal{B}$ 生成的开集族确实构成一个拓扑。
\end{proposition}

\subsection{子基}

\begin{definition}[子基]
设 $X$ 是集合，$\mathcal{S}\subseteq\mathcal{P}(X)$ 且 $\bigcup_{S\in\mathcal{S}}S=X$。由 $\mathcal{S}$ 的有限交作为基生成的拓扑，称为由 $\mathcal{S}$ 生成的拓扑；$\mathcal{S}$ 称为这个拓扑的一组\textbf{子基}。
\end{definition}

\begin{example}
在 $\mathbb{R}$ 上，所有形如 $(-\infty,a)$ 和 $(a,\infty)$ 的集合构成通常拓扑的一组子基。有限交可以给出开区间 $(a,b)$。
\end{example}

\subsection{可数性公理}

\begin{definition}[第一可数与第二可数]
拓扑空间 $X$ 称为\textbf{第一可数}，如果每个点都有一组可数的邻域基。$X$ 称为\textbf{第二可数}，如果它的拓扑有一组可数基。
\end{definition}

\begin{example}
$\mathbb{R}^n$ 是第二可数的。所有以有理点为中心、以正有理数为半径的开球构成一组可数基。
\end{example}

\begin{proposition}
第二可数空间一定是第一可数空间。
\end{proposition}

\section{闭集、闭包、内部与边界}

\begin{definition}[闭集]
设 $(X,\mathcal{T})$ 是拓扑空间。子集 $F\subseteq X$ 称为闭集，如果 $X\setminus F$ 是开集。
\end{definition}

\begin{proposition}[闭集的性质]
闭集满足以下性质：
\begin{enumerate}
    \item $\varnothing$ 和 $X$ 都是闭集；
    \item 任意多个闭集的交仍为闭集；
    \item 有限多个闭集的并仍为闭集。
\end{enumerate}
\end{proposition}

这些性质与开集性质互为对偶，本质上来自德摩根律。

\subsection{闭包与内部}

\begin{definition}[闭包]
设 $A\subseteq X$。包含 $A$ 的所有闭集的交称为 $A$ 的\textbf{闭包}，记作 $\overline{A}$。
\end{definition}

\begin{definition}[内部]
包含在 $A$ 中的所有开集的并称为 $A$ 的\textbf{内部}，记作 $A^\circ$ 或 $\operatorname{Int}(A)$。
\end{definition}

\begin{definition}[边界]
集合 $A$ 的\textbf{边界}定义为
\[
    \partial A=\overline{A}\setminus A^\circ.
\]
等价地，$x\in\partial A$ 当且仅当 $x$ 的每个邻域都同时与 $A$ 和 $X\setminus A$ 相交。
\end{definition}

\begin{example}
在 $\mathbb{R}$ 的通常拓扑中，若 $A=(0,1)$，则
\[
    \overline{A}=[0,1],\qquad A^\circ=(0,1),\qquad \partial A=\{0,1\}.
\]
若 $A=\mathbb{Q}$，则
\[
    \overline{\mathbb{Q}}=\mathbb{R},\qquad \mathbb{Q}^\circ=\varnothing.
\]
\end{example}

\begin{proposition}[用开邻域刻画闭包]
点 $x\in\overline{A}$ 当且仅当 $x$ 的每个开邻域都与 $A$ 相交。
\end{proposition}

\subsection{聚点与稠密性}

\begin{definition}[聚点]
设 $A\subseteq X$。点 $x\in X$ 称为 $A$ 的\textbf{聚点}，如果 $x$ 的每个邻域都含有某个不同于 $x$ 的 $A$ 中点。
\end{definition}

\begin{definition}[稠密子集]
如果 $\overline{A}=X$，则称 $A$ 在 $X$ 中\textbf{稠密}。
\end{definition}

\begin{example}
$\mathbb{Q}$ 和 $\mathbb{R}\setminus\mathbb{Q}$ 都在 $\mathbb{R}$ 中稠密。这个事实说明，稠密性不是“占据大多数点”的意思，而是指每个非空开集都能碰到这个集合。
\end{example}

\section{构造新的空间}

\subsection{子空间拓扑}

\begin{definition}[子空间拓扑]
设 $(X,\mathcal{T})$ 是拓扑空间，$Y\subseteq X$。定义
\[
    \mathcal{T}_Y=\{U\cap Y:U\in\mathcal{T}\}.
\]
则 $\mathcal{T}_Y$ 是 $Y$ 上的拓扑，称为从 $X$ 继承来的\textbf{子空间拓扑}。
\end{definition}

\begin{example}
区间 $[0,1]$ 作为 $\mathbb{R}$ 的子空间时，集合 $[0,1/2)$ 在 $[0,1]$ 中是开集，因为
\[
    [0,1/2)=(-1/2,1/2)\cap[0,1].
\]
但它不是 $\mathbb{R}$ 中的开集。
\end{example}

\subsection{积拓扑}

\begin{definition}[积拓扑]
设 $X,Y$ 是拓扑空间。$X\times Y$ 上的积拓扑以所有
\[
    U\times V,
\]
为基，其中 $U\subseteq X$ 开，$V\subseteq Y$ 开。
\end{definition}

\begin{example}
$\mathbb{R}^2$ 的通常拓扑等于 $\mathbb{R}\times\mathbb{R}$ 的积拓扑。
\end{example}

\newpage

\begin{proposition}[映入积空间的连续性]
映射 $f:Z\to X\times Y$ 连续，当且仅当它的两个分量映射
\[
    \pi_X\circ f:Z\to X,
    \qquad
    \pi_Y\circ f:Z\to Y
\]
都连续。
\end{proposition}

\subsection{商拓扑}

\begin{definition}[商拓扑]
设 $X$ 是拓扑空间，$\sim$ 是 $X$ 上的等价关系，并设 $q:X\to X/\sim$ 是自然投影。集合 $X/\sim$ 上的商拓扑定义为：$U\subseteq X/\sim$ 开，当且仅当
\[
    q^{-1}(U)
\]
在 $X$ 中开。
\end{definition}

\begin{example}[圆作为商空间]
把闭区间 $[0,1]$ 的两个端点粘合，也就是规定 $0\sim1$，得到的商空间同胚于圆 $S^1$。
\end{example}

\begin{remark}
商拓扑形式化了“粘合空间”的想法。圆、环面、莫比乌斯带等空间都可以通过适当的粘合构造出来。
\end{remark}

\section{分离公理}

拓扑空间中的点未必能被开集很好地区分。分离公理用于衡量这种区分能力。

\begin{definition}[$T_0$、$T_1$ 与 Hausdorff 空间]
设 $X$ 是拓扑空间。
\begin{enumerate}
    \item $X$ 称为 $T_0$ 空间，如果任意两个不同点中，至少有一个点有开邻域不包含另一个点；
    \item $X$ 称为 $T_1$ 空间，如果任意两个不同点 $x,y$，都存在开集 $U,V$，使得 $x\in U,y\notin U$ 且 $y\in V,x\notin V$；
    \item $X$ 称为 Hausdorff 空间或 $T_2$ 空间，如果任意两个不同点 $x,y$，存在不相交开集 $U,V$，使得 $x\in U$ 且 $y\in V$。
\end{enumerate}
\end{definition}

\begin{proposition}
Hausdorff 空间一定是 $T_1$ 空间，$T_1$ 空间一定是 $T_0$ 空间。
\end{proposition}

\begin{example}
任意度量空间都是 Hausdorff 空间。若 $x\ne y$，取
\[
    r=\frac12 d(x,y),
\]
则 $B_r(x)$ 和 $B_r(y)$ 不相交。
\end{example}

\begin{proposition}
拓扑空间是 $T_1$ 空间，当且仅当每个单点集 $\{x\}$ 都是闭集。
\end{proposition}

\begin{theorem}[Hausdorff 空间中极限唯一]
在 Hausdorff 空间中，若一个网或序列的极限存在，则极限唯一。特别地，在度量空间中，收敛序列的极限唯一。
\end{theorem}

\section{连续性与同胚}

\subsection{连续映射}

\begin{definition}[连续性]
设 $X,Y$ 是拓扑空间。映射 $f:X\to Y$ 称为连续，如果对任意开集 $V\subseteq Y$，原像 $f^{-1}(V)$ 都是 $X$ 中的开集。
\end{definition}

\begin{proposition}[闭集形式]
映射 $f:X\to Y$ 连续，当且仅当对任意闭集 $F\subseteq Y$，原像 $f^{-1}(F)$ 都是 $X$ 中的闭集。
\end{proposition}

\begin{proposition}[连续映射的复合]
若 $f:X\to Y$ 和 $g:Y\to Z$ 连续，则 $g\circ f:X\to Z$ 连续。
\end{proposition}

\subsection{同胚}

\begin{definition}[同胚]
双射 $f:X\to Y$ 称为\textbf{同胚}，如果 $f$ 连续且反函数 $f^{-1}:Y\to X$ 也连续。若存在同胚 $X\to Y$，则称 $X$ 与 $Y$ 同胚，记作
\[
    X\cong Y.
\]
\end{definition}

同胚是拓扑学中的“相同”。两个同胚空间可能在几何形状上看起来不同，但它们的开集结构完全等价，因此具有相同的拓扑性质。

\begin{example}
开区间 $(0,1)$ 与 $\mathbb{R}$ 同胚。例如映射
\[
    f(x)=\tan\left(\pi x-\frac{\pi}{2}\right)
\]
给出一个同胚 $(0,1)\to\mathbb{R}$。
\end{example}

\begin{remark}
连续双射不一定是同胚，因为反函数可能不连续。因此，在拓扑学中“双射”本身不足以表示两个空间相同。
\end{remark}

\section{连通性}

\subsection{定义}

\begin{definition}[分离]
拓扑空间 $X$ 的一个\textbf{分离}是两个非空、不相交开集 $U,V$，使得
\[
    X=U\cup V.
\]
\end{definition}

\begin{definition}[连通空间]
如果拓扑空间 $X$ 不存在分离，则称 $X$ 是\textbf{连通的}。
\end{definition}

\begin{proposition}[等价刻画]
拓扑空间 $X$ 连通，当且仅当 $X$ 中既开又闭的子集只有 $\varnothing$ 和 $X$ 本身。
\end{proposition}

\begin{theorem}[区间的连通性]
$\mathbb{R}$ 中的任意区间都是连通的。
\end{theorem}

\subsection{连通性的连续像}

\begin{theorem}
连通空间在连续映射下的像仍然连通。也就是说，如果 $f:X\to Y$ 连续且 $X$ 连通，则 $f(X)$ 连通。
\end{theorem}

\begin{theorem}[介值定理]
设 $X$ 连通，$f:X\to\mathbb{R}$ 连续。如果 $a,b\in f(X)$ 且 $a<c<b$，则 $c\in f(X)$。
\end{theorem}

这个定理说明，传统微积分中的介值定理本质上是连通性的结果。

\subsection{道路连通}

\begin{definition}[道路连通]
拓扑空间 $X$ 称为\textbf{道路连通}，如果对任意 $x,y\in X$，都存在连续映射
\[
    \gamma:[0,1]\to X
\]
使得 $\gamma(0)=x$ 且 $\gamma(1)=y$。
\end{definition}

\begin{proposition}
道路连通空间一定连通。
\end{proposition}

\begin{remark}
连通空间不一定道路连通。经典例子是拓扑学家的正弦曲线及其闭包。这说明，连通性只阻止空间被开集拆成两块，而道路连通还要求任意两点之间存在连续路径。
\end{remark}

\section{紧致性}

\subsection{开覆盖}

\begin{definition}[开覆盖]
设 $A\subseteq X$。一族开集 $\{U_\alpha\}_{\alpha\in I}$ 称为 $A$ 的\textbf{开覆盖}，如果
\[
    A\subseteq\bigcup_{\alpha\in I}U_\alpha.
\]
若存在有限多个指标 $\alpha_1,\dots,\alpha_n$ 使得
\[
    A\subseteq U_{\alpha_1}\cup\cdots\cup U_{\alpha_n},
\]
则称这是一组有限子覆盖。
\end{definition}

\begin{definition}[紧致性]
拓扑空间 $X$ 称为\textbf{紧致}，如果 $X$ 的每个开覆盖都有有限子覆盖。子集 $K\subseteq X$ 称为紧致，如果它作为子空间是紧致的。
\end{definition}

\begin{example}
有限拓扑空间总是紧致的，因为任意开覆盖只需要从每个点所在的开集中选出有限多个即可。
\end{example}

\begin{example}
开区间 $(0,1)$ 在通常拓扑下不是紧致的。开覆盖
\[
    \left\{\left(\frac1n,1\right):n\ge2\right\}
\]
覆盖 $(0,1)$，但没有有限子覆盖。
\end{example}

\subsection{\texorpdfstring{$\mathbb{R}^n$}{Rn} 中的紧致性}

\begin{theorem}[Heine--Borel 定理]
$\mathbb{R}^n$ 中的子集紧致，当且仅当它闭且有界。
\end{theorem}

\begin{remark}
闭且有界推出紧致，是欧氏空间的特殊性质。在一般度量空间或拓扑空间中，闭且有界未必意味着紧致。
\end{remark}

\begin{theorem}[Bolzano--Weierstrass 定理]
$\mathbb{R}^n$ 中任意有界序列都有收敛子列。
\end{theorem}

\begin{remark}
在欧氏空间中，紧致性、序列紧致性以及有界闭集之间有非常紧密的联系；但在一般拓扑空间中，这些概念可能分离。
\end{remark}

\subsection{紧致性的连续像}

\begin{theorem}
紧致空间在连续映射下的像仍然紧致。
\end{theorem}

\begin{theorem}[极值定理]
若 $K$ 紧致，$f:K\to\mathbb{R}$ 连续，则 $f$ 在 $K$ 上取得最大值和最小值。
\end{theorem}

\begin{theorem}[从紧致到 Hausdorff]
若 $X$ 紧致，$Y$ Hausdorff，且 $f:X\to Y$ 是连续双射，则 $f$ 是同胚。
\end{theorem}

\newpage

\section{重新看度量空间}

\begin{definition}[可度量化空间]
拓扑空间 $(X,\mathcal{T})$ 称为\textbf{可度量化}，如果存在度量 $d$，使得 $d$ 诱导出的拓扑恰好等于 $\mathcal{T}$。
\end{definition}

\begin{example}
离散拓扑总是可度量化的。定义
\[
    d(x,y)=
    \begin{cases}
    0,&x=y,\\
    1,&x\ne y,
    \end{cases}
\]
则由该度量诱导的拓扑就是离散拓扑。
\end{example}

\begin{remark}
不是每个拓扑空间都可度量化。度量空间具有许多额外性质，例如 Hausdorff 性、第一可数性等。可度量化定理研究的是哪些拓扑条件足以保证一个拓扑来自某个度量。
\end{remark}

\section{Hausdorff 条件之外的分离公理}

公理 $T_0$、$T_1$ 与 $T_2$ 比较点的邻域。进一步的分离公理则先比较点与闭集，再比较两个闭集。这种递进并非只是形式上的安排。分析学中的许多构造都需要用连续函数区分不同集合，而仅凭 Hausdorff 条件无法保证这类函数的存在。

不同文献采用的约定并不完全一致。在本书中，除非特别说明，\emph{正则空间}和\emph{正规空间}并不预先假定满足 $T_1$。因此，$T_3$ 表示正则且满足 $T_1$，$T_4$ 表示正规且满足 $T_1$。

\begin{definition}[正则性与正规性]
设 $X$ 是拓扑空间。
\begin{enumerate}
    \item 如果对任意 $x\in X$ 和不含 $x$ 的闭集 $F\subseteq X$，都存在不相交开集 $U,V$，使得 $x\in U$ 且 $F\subseteq V$，则称 $X$ 是\textbf{正则的}；
    \item 如果对任意闭集 $F$ 和任意 $x\notin F$，都存在连续函数 $f:X\to[0,1]$，满足 $f(x)=0$ 且 $f|_F=1$，则称 $X$ 是\textbf{完全正则的}；
    \item 如果对任意两个不相交闭集 $A,B\subseteq X$，都存在不相交开集 $U,V$，使得 $A\subseteq U$ 且 $B\subseteq V$，则称 $X$ 是\textbf{正规的}。
\end{enumerate}
\end{definition}

\begin{definition}[分离公理的层级]
空间按下表命名：
\[
\begin{array}{c|l}
T_0 & \text{Kolmogorov 空间},\\
T_1 & \text{所有单点集均为闭集},\\
T_2 & \text{Hausdorff 空间},\\
T_3 & \text{正则且满足 }T_1,\\
T_{3\frac12} & \text{完全正则且满足 }T_1,\\
T_4 & \text{正规且满足 }T_1.
\end{array}
\]
$T_{3\frac12}$ 空间也称为 \textbf{Tychonoff 空间}。
\end{definition}

\begin{proposition}[正则性的闭包刻画]
空间 $X$ 正则，当且仅当对每个点 $x$ 及每个包含 $x$ 的开集 $U$，都存在开集 $V$，使得
\[
    x\in V\subseteq \overline{V}\subseteq U.
\]
\end{proposition}

\begin{proof}
先设 $X$ 正则。集合 $F=X\setminus U$ 是不含 $x$ 的闭集。取不相交开集 $V,W$，使得 $x\in V$ 且 $F\subseteq W$。由于 $V\cap W=\emptyset$，$W$ 中没有点属于 $\overline{V}$。因此 $\overline{V}\subseteq X\setminus W\subseteq U$。

反过来，设上述闭包条件成立。令 $F$ 为闭集且 $x\notin F$。对 $U=X\setminus F$ 应用该条件，取 $V$ 满足 $x\in V\subseteq\overline{V}\subseteq U$。于是 $V$ 与 $X\setminus\overline{V}$ 是两个不相交开集，分别包含 $x$ 与 $F$。
\end{proof}

\begin{theorem}[度量空间是正规空间]
每个度量空间都是正规空间。因此每个度量空间都是 $T_4$ 空间，进而也是 $T_{3\frac12}$、$T_3$、$T_2$、$T_1$ 与 $T_0$ 空间。
\end{theorem}

\begin{proof}
设 $A,B$ 是度量空间 $(X,d)$ 中两个不相交闭集。定义
\[
    d(x,A)=\inf_{a\in A}d(x,a),\qquad d(x,B)=\inf_{b\in B}d(x,b).
\]
这两个函数都是连续的。令
\[
    U=\{x:d(x,A)<d(x,B)\},\qquad
    V=\{x:d(x,B)<d(x,A)\}.
\]
$U,V$ 都是开集且彼此不交。如果 $a\in A$，则 $d(a,A)=0$；由于 $B$ 闭且 $a\notin B$，有 $d(a,B)>0$。因此 $A\subseteq U$，同理 $B\subseteq V$。
\end{proof}

\begin{example}[为什么假设不可省略]
无限集上的余有限拓扑满足 $T_1$，但不是 Hausdorff 的，因为任意两个非空开集都有交。因此蕴含 $T_1\Rightarrow T_2$ 不成立。Sierpinski 空间 $X=\{0,1\}$ 取拓扑 $\{\emptyset,\{1\},X\}$，它满足 $T_0$，但不满足 $T_1$。这些例子说明较低层级的分离公理确实彼此不同。
\end{example}

\begin{theorem}[Urysohn 引理]
设 $X$ 是正规空间，$A,B\subseteq X$ 是两个不相交闭集。则存在连续函数 $f:X\to[0,1]$，使得
\[
    f|_A=0,\qquad f|_B=1.
\]
\end{theorem}

\begin{proof}[构造概要]
证明反复利用正规性，在开集之间插入闭包。对每个二进有理数 $r\in[0,1]$，构造开集 $U_r$，使得
\[
    A\subseteq U_0,
    \qquad \overline{U_r}\subseteq U_s\quad(r<s),
    \qquad X\setminus U_1\supseteq B.
\]
定义
\[
    f(x)=\inf\{r:x\in U_r\},
\]
并约定空集的下确界为 $1$。嵌套关系推出
\[
    f^{-1}([0,a))=\bigcup_{r<a}U_r,
    \qquad
    f^{-1}((a,1])=\bigcup_{r>a}(X\setminus\overline{U_r}),
\]
所以相关原像都是开集。故 $f$ 连续，并具有所要求的边界取值。
\end{proof}

\begin{theorem}[Tietze 延拓定理]
设 $X$ 是正规空间，$A\subseteq X$ 是闭集，$f:A\to[-1,1]$ 连续。则存在连续函数 $F:X\to[-1,1]$，使得 $F|_A=f$。
\end{theorem}

\begin{remark}
Urysohn 引理用一个连续实值函数分离两个闭集；Tietze 定理则把闭子空间上已经定义的连续函数延拓到整个空间。后一个定理更深，但其证明本质上是反复应用前一个定理来逼近尚未消去的误差。这些结果解释了为什么正规性是构造实值连续函数时自然出现的拓扑条件。
\end{remark}

\section{分支、局部连通性与拓扑学家的正弦曲线}

连通性是一个整体条件：整个空间不能分成两个非空开部分。道路连通性更具构造性：任意两点都能由连续道路连接。局部连通性则要求每一点附近存在连通邻域。这些条件彼此都不应混淆。

\begin{definition}[连通分支与道路分支]
对 $x\in X$，包含 $x$ 的所有连通子集之并称为 $x$ 的\textbf{连通分支}；包含 $x$ 的所有道路连通子集之并称为 $x$ 的\textbf{道路分支}。
\end{definition}

\begin{proposition}
连通分支是闭集、两两不交，并且是极大的连通子集。道路分支两两不交，并且是极大的道路连通子集，但在任意拓扑空间中未必是闭集。
\end{proposition}

\begin{proof}
如果两个连通子集有交，则它们的并仍然连通。因此，用来定义连通分支的并是连通的，而两个不同连通分支不能相交。连通集的闭包仍然连通，所以极大性迫使每个连通分支包含其闭包，因而它是闭集。道路分支的结论来自类似事实：两个相交道路连通集的并仍然道路连通。
\end{proof}

\begin{definition}[局部连通性]
如果 $x$ 的每个邻域都包含一个含 $x$ 的连通开邻域，则称空间 $X$ 在 $x$ 处\textbf{局部连通}。如果 $x$ 的每个邻域都包含一个含 $x$ 的道路连通开邻域，则称 $X$ 在 $x$ 处\textbf{局部道路连通}。若相应条件在每一点成立，则称空间局部连通或局部道路连通。
\end{definition}

\begin{theorem}
如果 $X$ 局部道路连通，那么它的连通分支与道路分支重合，并且每个分支都是开集。
\end{theorem}

\begin{proof}
固定 $x\in X$，设 $P$ 是它的道路分支。对每个 $y\in P$，取 $y$ 的一个道路连通开邻域 $U_y$。$U_y$ 中每一点都能与 $y$ 连接，而 $y$ 又能与 $x$ 连接，因此 $U_y\subseteq P$。故 $P$ 是开集。道路分支构成一个由不相交开集组成的划分。连通子集不可能同时与该划分中的两个不同部分相交，所以 $x$ 的连通分支包含于 $P$。反向包含关系由道路连通集必为连通集立即得到。
\end{proof}

\begin{example}[拓扑学家的正弦曲线]
令
\[
    S=\left\{\left(x,\sin\frac1x\right):0<x\leq 1\right\}
      \cup\left(\{0\}\times[-1,1]\right)\subseteq\mathbb{R}^2.
\]
图像部分
\[
    G=\left\{\left(x,\sin\frac1x\right):0<x\leq 1\right\}
\]
是连通区间 $(0,1]$ 的连续像，因此 $G$ 连通。它的闭包恰好是 $S$，而连通集的闭包仍连通，所以 $S$ 连通。

但是，$S$ 中不存在一条道路把竖直线段上的点与 $G$ 中的点连接起来。假设这样的道路存在。当第一坐标沿正值趋于 $0$ 时，第二坐标必须跟随 $\sin(1/x)$，从而在接近 $-1$ 与 $1$ 的值之间无限振荡。在道路第一次离开竖直线段的时刻，连续性将不可能成立。因此 $S$ 连通但不道路连通，并且它在竖直线段上的各点处也不局部连通。
\end{example}

\begin{remark}
这个例子的价值在于：失败并不是由彼此相距很远的不连通部分造成的。图像在整条竖直线段上聚集，因此取闭包创造了连通性，却没有创造道路。它典型地说明，拓扑中的极限行为可能比几何直觉所暗示的更加微妙。
\end{remark}

\section{超越“闭且有界”的紧致性}

“紧致就是闭且有界”这一说法在有限维欧氏空间中成立，但它不是紧致性的定义，并且在许多其他空间中失效。开覆盖定义更加稳健，因为它在连续映射下保持不变，而且不依赖度量。若干等价刻画从不同角度揭示了紧致性如何把无穷多个局部条件转化为有限个整体条件。

\subsection{有限交性质}

\begin{definition}[有限交性质]
如果集合族 $\mathcal{F}$ 的每个有限子族都有非空交集，则称 $\mathcal{F}$ 具有\textbf{有限交性质}。
\end{definition}

\begin{theorem}[紧致性的闭集刻画]
拓扑空间 $X$ 紧致，当且仅当每个具有有限交性质的闭集族都有非空的全体交集。
\end{theorem}

\begin{proof}
设 $X$ 紧致，$\{F_i\}_{i\in I}$ 是全体交为空的闭集族。则 $\{X\setminus F_i\}_{i\in I}$ 是 $X$ 的开覆盖。若它有有限子覆盖，则相应有限多个 $F_i$ 的交为空，这与有限交性质矛盾。

反过来，设每个具有有限交性质的闭集族都有非空交。若某个开覆盖 $\{U_i\}_{i\in I}$ 没有有限子覆盖，则闭集 $X\setminus U_i$ 的任意有限交都非空。于是它们的全体交也应非空，这与 $U_i$ 覆盖 $X$ 矛盾。
\end{proof}

\subsection{度量空间中的列紧性}

\begin{definition}[列紧性]
如果每个序列都有一个收敛子列，并且其极限属于该空间，则称该空间是\textbf{列紧的}。
\end{definition}

\begin{theorem}
对度量空间而言，紧致性与列紧性等价。
\end{theorem}

\begin{proof}[证明思路]
若度量空间紧致，用有限多个半径为 $1$ 的球覆盖它。其中至少有一个球包含序列的无穷多项。再用有限多个半径为 $1/2$ 的球覆盖该球，并选择一个包含剩余项中无穷多项的球。如此归纳，便得到一列嵌套选取的项，它们之间的距离趋于 $0$。紧致性保证存在聚点，而选出的子列收敛到该点。

反过来，设空间列紧。首先，它必然全有界。否则存在 $\varepsilon>0$ 和序列 $(x_n)$，使得
\[
    d(x_m,x_n)\geq\varepsilon\qquad(m\neq n),
\]
而该序列不可能有收敛子列。

现在设 $\mathcal{U}$ 是一个开覆盖。我们断言 $\mathcal{U}$ 存在 Lebesgue 数。否则，对每个 $n$ 可选取 $x_n$，使得 $B_{1/n}(x_n)$ 不包含于 $\mathcal{U}$ 的任何一个成员。由列紧性，存在子列 $x_{n_k}\to x$。取 $U\in\mathcal{U}$ 和 $r>0$，使得 $B_r(x)\subseteq U$。当 $k$ 充分大时，同时有 $d(x_{n_k},x)<r/2$ 和 $1/n_k<r/2$，于是
\[
    B_{1/n_k}(x_{n_k})\subseteq B_r(x)\subseteq U,
\]
矛盾。最后，全有界性给出有限多个半径小于 Lebesgue 数的球，而每个球都包含于 $\mathcal{U}$ 的某个成员。这有限多个成员便构成有限子覆盖。
\end{proof}

\begin{theorem}[Lebesgue 数引理]
设 $(X,d)$ 是紧度量空间，$\mathcal{U}$ 是 $X$ 的开覆盖。则存在 $\delta>0$，使得每个直径小于 $\delta$ 的子集都包含于 $\mathcal{U}$ 的某个成员。
\end{theorem}

\begin{proof}
对每个 $x\in X$，取包含 $x$ 的 $U_x\in\mathcal{U}$，并取 $r_x>0$，使得 $B_{2r_x}(x)\subseteq U_x$。球族 $B_{r_x}(x)$ 覆盖 $X$，所以其中有限多个已经足够，记为 $B_{r_1}(x_1),\dots,B_{r_n}(x_n)$。令 $\delta=\min r_i$。若集合 $A$ 的直径小于 $\delta$，并取 $a\in A$，选取 $i$ 使得 $a\in B_{r_i}(x_i)$。则对每个 $y\in A$，
\[
    d(y,x_i)\leq d(y,a)+d(a,x_i)<\delta+r_i\leq 2r_i,
\]
所以 $A\subseteq B_{2r_i}(x_i)\subseteq U_{x_i}$。
\end{proof}

\begin{corollary}[紧致空间上的一致连续性]
从紧度量空间到度量空间的每个连续映射都是一致连续的。
\end{corollary}

\begin{proof}
设 $f:X\to Y$ 连续，固定 $\varepsilon>0$。对每个 $x\in X$，选取开邻域 $U_x$，使得 $f(U_x)$ 包含于以 $f(x)$ 为中心、半径为 $\varepsilon/2$ 的球。覆盖 $\{U_x\}$ 的一个 Lebesgue 数 $\delta$ 保证：只要 $d(p,q)<\delta$，二点集 $\{p,q\}$ 就包含于某个 $U_x$。因此
\[
    d_Y(f(p),f(q))<\varepsilon.
\]
\end{proof}

\subsection{有限乘积与 Tychonoff 定理的初等形式}

\begin{theorem}[有限乘积定理]
如果 $X$ 与 $Y$ 都紧致，那么 $X\times Y$ 在积拓扑下紧致。因此，紧致空间的任意有限乘积仍然紧致。
\end{theorem}

\begin{proof}[利用管状引理的证明]
设 $\mathcal{U}$ 是 $X\times Y$ 的一个开覆盖。固定 $x\in X$。切片 $\{x\}\times Y$ 与 $Y$ 同胚，因而紧致，所以 $\mathcal{U}$ 中有限多个成员即可覆盖该切片。它们的并是一个包含 $\{x\}\times Y$ 的开集。由管状引理，存在 $x$ 的开邻域 $N_x$，使得
\[
    N_x\times Y
\]
包含于上述有限并中。集合 $N_x$ 覆盖 $X$。由 $X$ 的紧致性，其中有限多个已经覆盖 $X$；收集与它们对应的 $\mathcal{U}$ 中有限子族，即得 $X\times Y$ 的有限覆盖。
\end{proof}

\begin{lemma}[管状引理]
设 $Y$ 紧致，$W\subseteq X\times Y$ 是开集，且 $\{x_0\}\times Y\subseteq W$。则存在 $x_0$ 的开邻域 $N$，使得 $N\times Y\subseteq W$。
\end{lemma}

\begin{proof}
对每个 $y\in Y$，选取基本开集 $N_y\subseteq X$ 和 $V_y\subseteq Y$，使得
\[
    (x_0,y)\in N_y\times V_y\subseteq W.
\]
集合 $V_y$ 覆盖 $Y$，故有限多个 $V_{y_1},\dots,V_{y_n}$ 已经足够。于是
\[
    N=N_{y_1}\cap\cdots\cap N_{y_n}
\]
是 $x_0$ 的开邻域，并且 $N\times Y\subseteq W$。
\end{proof}

\begin{remark}
完整的 Tychonoff 定理断言：紧致空间的任意乘积都是紧致的。其标准证明使用选择公理，通常借助超滤子或具有有限交性质的极大家族。上面的有限乘积定理包含了其中的几何核心，并足以应对大多数初等应用。
\end{remark}

\section{作为泛构造的积与商}

积拓扑与商拓扑的定义看起来并不对称：积拓扑由投影映射的原像生成，而商拓扑则通过一个满射的原像来检验。它们的共同特征是泛性质：每种拓扑都恰好使某类连续性问题变得尽可能容易。

\begin{theorem}[积的泛性质]
设 $Z$ 是拓扑空间，$f:Z\to X\times Y$，写成
\[
    f(z)=(f_1(z),f_2(z)).
\]
则 $f$ 连续，当且仅当两个坐标映射 $f_1:Z\to X$ 与 $f_2:Z\to Y$ 都连续。
\end{theorem}

\begin{proof}
如果 $f$ 连续，则 $f_1=\pi_X\circ f$ 与 $f_2=\pi_Y\circ f$ 都连续。反过来，对每个基本开集 $U\times V$，有
\[
    f^{-1}(U\times V)=f_1^{-1}(U)\cap f_2^{-1}(V),
\]
当 $f_1,f_2$ 连续时，它是开集。
\end{proof}

\begin{theorem}[商的泛性质]
设 $q:X\to Y$ 是商映射，$g:Y\to Z$ 是函数。则 $g$ 连续，当且仅当 $g\circ q:X\to Z$ 连续。
\end{theorem}

\begin{proof}
一个方向由连续映射的复合立即得到。反过来，设 $g\circ q$ 连续。对每个开集 $W\subseteq Z$，
\[
    q^{-1}(g^{-1}(W))=(g\circ q)^{-1}(W)
\]
在 $X$ 中是开集。由商拓扑的定义，$g^{-1}(W)$ 在 $Y$ 中是开集。
\end{proof}

\begin{example}[由区间得到圆周]
定义 $q:[0,1]\to S^1$ 为
\[
    q(t)=e^{2\pi i t}.
\]
该映射连续、满射，并且恰好把端点 $0$ 与 $1$ 识别。由于 $[0,1]$ 紧致而 $S^1$ 是 Hausdorff 空间，诱导出的双射
\[
    [0,1]/(0\sim1)\longrightarrow S^1
\]
是同胚。这一论证比图形直观更强：它证明了该商拓扑恰好就是圆周上的通常拓扑。
\end{example}

\begin{example}[一个非 Hausdorff 商空间]
取两份实直线，并把除原点外的对应点逐一识别，所得空间称为“双原点直线”。它在局部与 $\mathbb{R}$ 无法区分，但两个原点不能被不相交开集分离。因此，取商可能破坏 Hausdorff 性。每当通过粘合构造空间时，都必须检验分离公理，而不能默认它们成立。
\end{example}


\section{结语}

拓扑学把“接近”“连续”“边界”“整体连通”等观念从坐标和距离中解放出来。它提供了一种高度抽象但非常稳定的语言：只要知道开集，就能重建大量分析学中熟悉的概念。后续的复分析、微分几何、泛函分析乃至现代物理中的空间理论，都离不开这种拓扑语言。

\textbf{关键词：} 拓扑空间，开集，基，闭包，连续映射，同胚，连通性，紧致性。
